\documentclass[11pt]{article}

\usepackage[margin=1in]{geometry}
\usepackage{amsmath, amssymb, amsthm}
\usepackage{xcolor}
\usepackage{listings}
\lstset{
  basicstyle=\ttfamily\small,
  numbers=left,
  numberstyle=\tiny\color{gray},
  numbersep=8pt,
  stepnumber=1,
  frame=single,
  framesep=6pt,
  breaklines=true,
  columns=fullflexible,
  keepspaces=true,
  showstringspaces=false,
  xleftmargin=2.5em,
}
\usepackage{hyperref}

\hypersetup{colorlinks=true, linkcolor=black, urlcolor=black}

\newtheorem{theorem}{Theorem}
\newtheorem{definition}{Definition}

\newcommand{\code}[1]{\texttt{#1}}
\newcommand{\R}{\mathbb{R}}
\DeclareMathOperator*{\supp}{sup'}
\DeclareMathOperator*{\infp}{inf'}

% Minimal placeholder for the student to type over
\newcommand{\blank}{\par\vspace{2pt}\noindent\textbf{[YOUR ANSWER]}\par\vspace{6pt}}
\newcommand{\blankone}{\textbf{[YOUR ANSWER]}}

\title{Negation Flips \texorpdfstring{$\max$}{max} to \texorpdfstring{$\min$}{min}:\\
  Fard \& Pineau (2011), (20)$\to$(21) --- \emph{Student Fill-In Version}}
\author{}
\date{}

\begin{document}
\maketitle

\noindent\small\textit{Instructions: the structure, equations, and step labels
are given. Replace each \textbf{[YOUR ANSWER]} with your own explanation. If you
can fill every one correctly, you understand the proof.}\normalsize

\vspace{1em}

% ==========================================================================
\section{The Original Result}

\emph{What is the evaluation MDP $M'$, and how are $A'$ and $R'$ defined in terms
of the original?}
\blank

\begin{theorem}[Fard \& Pineau, Thm.~1]
\[
  Q^{\Pi}_M(s,a) \;=\; -\,Q^{*}_{M'}(s,a).
\]
\end{theorem}

\emph{In one line, what does the proof do to get from the Bellman equation on
$M'$ to the claim?}
\blank

\begin{align}
  Q^{*}_{M'}(s,a)
    &= \bar R'(s,a) + \gamma \sum_{s'} T(s,a,s') \max_{a' \in A'} Q^{*}_{M'}(s',a')
      \tag{19}\\[2pt]
  \Rightarrow\; -Q^{*}_{M'}(s,a)
    &= -\bar R'(s,a) - \gamma \sum_{s'} T(s,a,s') \max_{a' \in A'} Q^{*}_{M'}(s',a')
      \tag{20}\\[2pt]
  \Rightarrow\; -Q^{*}_{M'}(s,a)
    &= \bar R(s,a) + \gamma \sum_{s'} T(s,a,s') \min_{a' \in \Pi(s')}\!\bigl(-Q^{*}_{M'}(s',a')\bigr).
      \tag{21}
\end{align}

\emph{Between (20) and (21), which changes are purely definitional substitutions,
and which single change is the real inference?}
\blank

\emph{State the one substantive identity being used:}
\[
  \phantom{-\max_{a'} Q \;=\; \min_{a'}(-Q)}
\]
\blankone

\emph{What hypothesis (from Definition~1) is implicit here but cannot be skipped
in a formalization? Why is it needed?}
\blank

% ==========================================================================
\section{The Formalization, Step by Step}

\emph{How does the bare identity translate: what do the action set, the
$Q$-value, and $\max/\min$ become? What does the primed $\supp$/$\infp$ require
as an argument, and why?}
\blank

\begin{equation}
  -\,\supp_{a \in A} Q(a) \;=\; \infp_{a \in A}\bigl(-Q(a)\bigr).
  \tag{$\star$}
\end{equation}

\subsection*{Overall strategy}

\emph{To prove an equality of two reals, what do we prove instead, and how are
the two directions related?}
\blank

\textbf{Reshape.} \emph{What does a reshape step do, using which identities, and
does it change what is true?}
\blank

\textbf{Close.} \emph{What kind of fact closes a branch? List the four defining
facts about $\max/\min$ used.}
\blank

\subsection*{Direction 1: \texorpdfstring{$-\supp_a Q \le \infp_a(-Q)$}{-sup Q <= inf(-Q)}}

\begin{enumerate}
  \item \textbf{Reduce to elements.} \emph{Why does the RHS being a minimum let
    us reduce to a per-element statement?}
    \blank
    \[
      -\,\supp_{a'} Q(a') \;\le\; -Q(a). \qquad (\forall\, a \in A)
    \]
  \item \textbf{Fix an arbitrary $a$.} \emph{Why does proving it for a generic
    $a$ suffice for all elements? How is this different from a loop?}
    \blank
  \item \textbf{Reshape.} \emph{What is the shape of the goal, which identity
    applies, and what does the goal become?}
    \blank
    \[
      Q(a) \;\le\; \supp_{a'} Q(a').
    \]
  \item \textbf{Close.} \emph{Which defining fact finishes this, and why does it
    match?}
    \blank
\end{enumerate}

\subsection*{Direction 2: \texorpdfstring{$\infp_a(-Q) \le -\supp_a Q$}{inf(-Q) <= -sup Q}}

\emph{How does this direction relate to Direction 1?}
\blankone

\begin{enumerate}
  \item \textbf{Reshape \#1.} \emph{What shape is the goal, which identity
    applies, and what gets exposed?}
    \blank
    \[
      \supp_{a'} Q(a') \;\le\; -\,\infp_{a'}\bigl(-Q(a')\bigr).
    \]
  \item \textbf{Reduce to elements.} \emph{Why can we now reduce to a per-element
    statement?}
    \blank
    \[
      Q(a) \;\le\; -\,\infp_{a'}\bigl(-Q(a')\bigr). \qquad (\forall\, a \in A)
    \]
  \item \textbf{Fix an arbitrary $a$.}
  \item \textbf{Reshape \#2.} \emph{Same identity as \#1 --- what does it expose
    this time?}
    \blank
    \[
      \infp_{a'}\bigl(-Q(a')\bigr) \;\le\; -Q(a).
    \]
  \item \textbf{Close.} \emph{Which defining fact finishes this, and against
    which element?}
    \blank
\end{enumerate}

\emph{With both directions done, what lets us conclude the equality~$(\star)$,
and how does $(\star)$ relate to the paper's (20)$\to$(21) step?}
\blank

\subsection*{Two remarks}

\begin{itemize}
  \item \textbf{``Fix an arbitrary $a$'' is not a loop.} \emph{Explain.}
    \blank
  \item \textbf{Non-emptiness is load-bearing.} \emph{Explain what breaks
    without it, and how the mechanization treats it.}
    \blank
\end{itemize}

% ==========================================================================
\appendix
\section{Lean 4 Source --- annotate each line yourself}

\emph{For each tactic below, note: is it a} Reshape \emph{or a}
Close\emph{, and what is the goal before/after? (Full source intentionally given;
the comments are removed so you supply them.)}

\begin{lstlisting}
import Mathlib
open Finset

theorem neg_max_eq_min_neg
    {Action : Type*}
    (actions : Finset Action)
    (h_nonempty : actions.Nonempty)
    (Q : Action -> R) :
    - actions.sup' h_nonempty Q
      = actions.inf' h_nonempty (fun x => - Q x) := by
  apply le_antisymm
  . apply Finset.le_inf'
    intro a ha
    rw [neg_le_neg_iff]
    exact Finset.le_sup' Q ha
  . rw [le_neg]
    apply Finset.sup'_le
    intro a ha
    rw [le_neg]
    exact Finset.inf'_le (fun x => - Q x) ha
\end{lstlisting}

\emph{Questions to answer for the appendix:}
\begin{itemize}
  \item Why does \code{apply le\_antisymm} produce exactly two goals?
    \blankone
  \item What does \code{intro a ha} destructure, and where does the type of
    \code{a} come from? \blankone
  \item Why is \code{h\_nonempty} passed around even though it contributes no
    number to the answer? \blankone
  \item In the last line, what are the two arguments to \code{Finset.inf'\_le},
    and what does the applied lambda reduce to? \blankone
\end{itemize}

\end{document}
